% !TeX root = ./proposta-pytorch.tex
\documentclass[a4paper,12pt, brazil]{article}

% --- PREÂMBULO ---
% 1. Carregue os pacotes/dependências do arquivo
% █▀█ █▀█ █▀▀ ▄▀█ █▀▄▀█ █▄▄ █   █▀▀
% █▀▀ █▀▄ ██▄ █▀█ █ ▀ █ █▄█ █▄▄ ██▄
%::::::::::::::::::::::::::::::::::::::::::::::::::::::
%                      CORE SETUP                      
%                Document-Wide Settings                
%......................................................
\usepackage[utf8]{inputenc}	         % Input encoding: Assumes UTF-8 text editor
\usepackage[T1]{fontenc}             % Output enconding: 8-bit, for better glyph support
\usepackage[margin=2.5cm]{geometry}  % Set page margins
\usepackage[brazil]{babel}           % Language-specific settings

%::::::::::::::::::::::::::::::::::::::::::::::::::::::
%                       GRAPHICS                       
%                Generate/Import Images                 
%......................................................
%% NATIVE PROGRAMMATIC DIAGRAMS & PLOTS
%\usepackage{tikz}		        % Create vector graphics programmatically
%\usetikzlibrary{arrows.meta}	% Additional arrow styles
%\usepackage{pgfplots}		    % Tool for creating function plots and charts
%\pgfplotsset{compat=1.18}	    % Enforce version 1.18 for consistency

% EXTERNAL IMAGE MANAGEMENT
\usepackage{graphicx}		    % LaTeX standard for including external image files

%::::::::::::::::::::::::::::::::::::::::::::::::::::::
%                    MATH PACKAGES                     
%......................................................
% AMS (American Mathematical Society)
\usepackage{amssymb} 	% AMS's package for mathematical symbols
\usepackage{amsmath}	% AMS's package for mathematical typesetting
\usepackage{mathtools}  % Extends amsmath's functionalities (i.e.: a superset)
\usepackage{amsthm}    	% AMS's package for creating theorem-like environments


% LOADED LAST TO AVOID DEPENDENCY ERRORS
%::::::::::::::::::::::::::::::::::::::::::::::::::::::
%                     BIBLIOGRAPHY                      
%                   Manage References                   
%......................................................
\usepackage[
	backend=biber,	% Use Biber backend (more powerful than BibTex)
	style=ieee,     %
	sorting=nyt     % 'nyt' stands for name, year, title
]{biblatex}

%::::::::::::::::::::::::::::::::::::::::::::::::::::::
%                     FINALIZATION                      
%    Packages That Patch Others. MUST BE LOADED LAST    
%......................................................
\usepackage{hyperref}	% Add hyperlinks to references, ToC and URLs
\hypersetup{
	colorlinks=true,	% true: colored links; false: boxed links
	allcolors=NavyBlue,	% Set all link types to a single color
}

% 2. Configuração dos 'File Paths' (do inglês, 'Caminho de Arquivos')
\addbibresource{referencias.bib}
\graphicspath{ {./imagens/} }

% --- INÍCIO DO DOCUMENTO ---
\begin{document}
    % --- TÍTULO ---
    % !TeX root = ./proposta-pytorch.tex
\begin{center}
	\begin{figure}[h!]
		\centering
		\includegraphics[height=1.5cm]{logoICMC.png}
	\end{figure}

	\rule{0.95\textwidth}{1.5pt}

	\vspace*{2mm}

	{AEX \par}

	\vspace*{1mm}

	{\large Trabalho 1\par}

	\vspace*{3mm}

	{\Large\textbf{Pytorch do Zero}\par}
	
	\vspace*{3mm}

	{\large\textbf{\emph{Métodos da Bisseção, Newton-Rhapson e das Secantes}}}

	\vspace*{8mm}

	\begin{tabular}{rrcl}
        Docente Responsável: & Prof. Dr. Lucas Pascotti Valem \\
        Alunos: & \\
        & Amanda Araújo & -- & XXXXXXXXXX \\
		& Cody Stefano Barham Setti & -- & 4856322 \\
		& Ian Cavalcanti de Holanda Bezerra & -- & XXXXXXXXXX
	\end{tabular}

	\vspace*{8mm}
	
	{\normalsize \today\par}

	\vspace*{2mm}

	{\normalsize São Carlos, SP, Brasil\par}

	\vspace*{2mm}

	\rule{0.95\textwidth}{1.5pt}
\end{center}
\vspace*{5mm}

    \newpage
    
    % --- CORPO DO ARTIGO ---
    \section*{Proposta da AEX}
    
    \subsection*{Aulas}
    \subsubsection*{Bloco 1 -- IA do Zero}
    \begin{itemize}
        \item Vídeo 1 -- Introdução e Revisão
        \begin{itemize}
            \item Introdução ao curso;
            \item Machine Learning, introdução ao problema do pipeline, complexidade, informação;
            \item Environment, bibliotecas, colab. Configurando máquina.
        \end{itemize}
    
        \item Vídeo 2 -- Revisão de Conceitos Necessários
        \begin{itemize}
            \item Revisão Matemática: funções ($\mathbb{R}^n$ em $\mathbb{R}^n$), derivadas parciais, gradiente;
            \item Parte 1: Gradientes e Backprop.
        \end{itemize}
            
        \item Vídeo 3 -- \emph{PyTorch Autograd} do zero
            \begin{itemize}
                \item Parte 2: Model class, exemplo seno, primeiro MLP, e levando conceitos a Torch, torch optim;
                \item Parte 3: O que eh uma GPU, activation functions, revisão conceitos.
            \end{itemize} 
    \end{itemize}
    
    \subsubsection*{Bloco 2 -- Arquiteturas de Redes Neurais}
    \begin{itemize}
        \item Vídeo 4 -- Introdução às Arquiteturas de Redes Neurais
        \begin{itemize}
            \item Torch \verb|NN.module| e resolvendo MNIST (CPU, GPU);
            \item Problemas comuns -- Parte 1: \emph{Over-fitting}, \emph{exploding gradients} e \emph{vanishing gradients};
            \item Problemas comuns -- Parte 2: Layer Normalization, Batch Normalization, Shuffle correspondence.
        \end{itemize}
        
        \item Vídeo 5 -- Mais Arquiteturas de Redes Neurais e Prática
        \begin{itemize}
            \item CNNs;
            \item RNN.
        \end{itemize}
        
        \item Vídeo 6 -- Prática.
        \begin{itemize}
            \item Prática. \emph{Kaggle to Solved NN}.
        \end{itemize}
    \end{itemize}
    
    \subsubsection*{Bloco 3 -- Dados e Práticas de Produção}
    \begin{itemize}
        \item Vídeo 7 -- \emph{Embeddings} e \emph{Transfer Learning}.
        \begin{itemize}
            \item \emph{Embeddings}
            \item Modelos Pré-Treinados (LLMs, TTS, STT, Stable Diffusion);
            \item Pontos importantes de modelos grandes em GPU;
            \item \emph{Transfer-learning};
            \item \emph{Fine tunning}.
        \end{itemize}
        
        
        \item Vídeo 8 -- Dados: Extração, Contaminação, Criação (ChemAI)
    
        
        \item Vídeo 9 -- Infra-estrutura para IA
        
        \item Vídeo 10 -- Colocando IA em produção
    \end{itemize}
    
        
    \section*{Material}
    \begin{itemize}
        \item Bloco 1:
        \begin{itemize}
            \item Andrej Karpathy -- Neural Networks: Zero to Hero \cite{karpathy} (vídeos 1-2).
        \end{itemize}
        
        \item Bloco 2:
        \begin{itemize}
            \item MIT 6.034: Artificial Intelligence \cite{mit_ai};
            \item MIT 6.S191: Introduction to Deep Learning \cite{mit_deep_learning}.
        \end{itemize}
        
        \item Bloco 3:
        \begin{itemize}
            \item AI-Chem (material próprio);
            \item Andrej Karpathy -- Neural Networks: Zero to Hero \cite{karpathy} (vídeos 3-5);
            \item MIT 6.S191: Introduction to Deep Learning \cite{mit_deep_learning}.
        \end{itemize}
    
    \end{itemize}    
    
    \section*{Horas -- AEX}
    \begin{table}[h!]
        \centering
        \begin{tabular}{cc}
            \toprule
            \textbf{Atividade} & \textbf{Tempo (horas/aula/discente)} \\
            \midrule
            Preparo de material didático & $4$ \\
            Gravação de vídeo-aula & $1$ \\
            Edição de vídeo-aula & $3$ \\
            Atendimento síncrono & $2$ \\
            \bottomrule
        \end{tabular}
        \caption{}
        \label{tbl:horas-de-extensao}
    \end{table}
    \textbf{Total de vídeo-aulas: 10}

    \textbf{Total de horas}: $10\times(4\,\text{h} + 1\,\text{h} + 3\,\text{h} + 2\,\text{h}) = \boldsymbol{100}$ \textbf{horas/discente}.
    
    % --- BIBLIOGRAFIA ---
    \printbibliography
\end{document}
